\chapter{Message framing}

\section{The identifier test}\label{sec:identifier}
\estab\ The instrument accepts exactly two System Exclusive identifiers, and they are
literal immediate operands, not settings:

\begin{center}
\begin{tabular}{lll}
\toprule
identifier & meaning & instruction \\
\midrule
\bytes{0x50} & Matsushita (Technics) & \texttt{cp E,0x50} at \bytes{0xFA5850} \\
\bytes{0x7E} & Universal Non-Real Time & \texttt{cp E,0x7E} at \bytes{0xFA5855} \\
\bottomrule
\end{tabular}
\end{center}

There is no third comparison, no table lookup, and no read of any stored setting on this
path. \textbf{There is no device or unit number at this level} --- the byte after the
identifier is already treated as message body.

\section{How a message is tracked}
\estab\ \bytes{0xF0} is first handled as an ordinary status byte and stored as running
status. The SysEx logic runs when the \emph{next} byte arrives, so the byte the identifier
test inspects is message byte index~1.

A single RAM byte, \bytes{(0xA9)}, holds the whole receive state:

\begin{center}
\begin{tabular}{cll}
\toprule
bit & meaning & set at \\
\midrule
0 & armed --- an \bytes{F0} was seen & \bytes{0xFA584D} \\
1 & identifier accepted & \bytes{0xFA585A} \\
2 & a message completed cleanly (byte value \bytes{0x04}) & \bytes{0xFA54E2} \\
4 & aborted by a stray status byte (byte value \bytes{0x10}) & \bytes{0xFA54EE} \\
5 & suppress further SysEx & \bytes{0xFB6151} \\
\bottomrule
\end{tabular}
\end{center}

\estab\ A UART receive error clears the byte outright at \bytes{0xFA5426}.

\begin{tcolorbox}[colback=blue!4!white,colframe=blue!40!black,title=A correction to the
firmware's own annotations]
Both the block header above \sym{MIDI\_RX\_SysExStart} and the project's
\texttt{FINDINGS-midi-port.md} state that bit~5 is ``set nowhere in the converted code''.
It \emph{is} set, at \bytes{0xFB6151}, in the synchronous receiver on a SysEx framing
error. Bits~2 and~4 are written but no reader of them was found anywhere in
\texttt{prom\_a}.
\end{tcolorbox}

\section{What happens to another maker's message}
\estab\ An identifier that is neither \bytes{0x50} nor \bytes{0x7E} leaves bit~0 set and
bit~1 clear. Every following byte still reaches the data handler and is discarded there,
so the instrument tracks the foreign message to its \bytes{F7} without acting on it and
without raising any error. Outside a bulk-dump session this is silent.

\section{Termination, abort, and real time}
\estab\ Termination is decided in the status-byte path, not in the SysEx routines:

\begin{itemize}
\item \bytes{F7} with bit~5 clear forwards the \bytes{F7} and leaves \bytes{(0xA9)} =
  \bytes{0x04}, the ``complete message received'' residue.
\item Any other status byte in \bytes{0x80}--\bytes{0xF6} aborts: \bytes{(0xA9)} =
  \bytes{0x10} and nothing is forwarded.
\item System Real Time (\bytes{0xF8}--\bytes{0xFF}) is split off earlier, at
  \bytes{0xFA54B6}, so a clock or active-sensing byte \emph{inside} a System Exclusive
  message does not break it. This is correct MIDI 1.0 behaviour.
\end{itemize}

\estab\ Conversely, while a SysEx is armed or in progress the instrument ignores
\bytes{0xF8} \textsc{timing clock}, \bytes{0xFA} \textsc{start}, \bytes{0xFB}
\textsc{continue} and \bytes{0xFC} \textsc{stop} for transport purposes
(\bytes{0xFA5520}). The block sender likewise spins until an inbound message finishes
before transmitting (\bytes{0xFB99CC}).
